% Trustworthy AI terms, in alphabetical order.
\domain{Trustworthy AI}

\begin{entry}{accountability}{Accountability}
  \defn{hleg-2019}{A requirement for trustworthy AI: mechanisms are in place
    to ensure responsibility and accountability for AI systems and their
    outcomes, before and after development, deployment, and use. These
    include auditability, minimising and reporting negative impacts,
    trade-offs, and redress.}
  \defn[§3.4]{nist-ai-rmf}{Accountability presupposes transparency. It
    concerns who answers for an AI system's outcomes, and it depends on
    organisational practices and governing structures for reducing harms,
    such as risk management and clear roles and responsibilities.}
  \defn{bovens-2007}{In public governance, a relationship between an actor
    and a forum in which the actor must explain and justify its conduct, the
    forum can pose questions and pass judgement, and the actor may face
    consequences.}
  \note{Bovens's relational definition clarifies that accountability needs
    an identifiable party who must answer, and someone they must answer to.
    Transparency alone does not provide either.}
  \related{transparency,auditability,ai-governance}
\end{entry}

\begin{entry}{adversarial-example}{Adversarial Example}
  \defn{szegedy-2014}{An input formed by applying a small, often
    imperceptible, perturbation to a correctly classified example, chosen so
    that a neural network misclassifies it.}
  \defn{goodfellow-2015}{Inputs formed by intentionally adding worst-case
    perturbations. The authors argue that vulnerability to them comes mainly
    from the models' largely linear behaviour in high-dimensional spaces,
    not from non-linearity or overfitting.}
  \defn{nist-ai-100-2}{The product of an \emph{evasion attack}: at
    deployment time, an adversary modifies a test sample to change the
    model's prediction to one the attacker chooses.}
  \related{robustness,data-poisoning,security-resilience}
\end{entry}

\begin{entry}{ai-governance}{AI Governance}
  \defn{mantymaki-2022}{A system of rules, practices, processes, and
    technological tools used to ensure that an organisation's use of AI
    technologies is aligned with its strategies, objectives, and values,
    fulfils legal requirements, and meets principles of ethical AI.}
  \defn{iso-42001}{Realised through an \emph{AI management system}: the
    interrelated policies, objectives, and processes an organisation
    establishes for the responsible development, provision, or use of AI
    systems.}
  \defn[§5.1]{nist-ai-rmf}{The \textsc{Govern} function: a culture of risk
    management is cultivated and present across the organisation. It cuts
    across and informs the \textsc{Map}, \textsc{Measure}, and
    \textsc{Manage} functions.}
  \related{accountability,ai-risk,auditability}
\end{entry}

\begin{entry}{ai-risk}{AI Risk}
  \defn*[Art.~3(2)]{eu-ai-act}{The combination of the probability of an
    occurrence of harm and the severity of that harm.}
  \defn[§1.1]{nist-ai-rmf}{A composite measure of an event's probability of
    occurring and the magnitude or degree of the consequences. AI risks can
    affect individuals, groups, organisations, communities, society, the
    environment, and the planet.}
  \defn{iso-31000}{The effect of uncertainty on objectives, where an effect
    is a deviation from the expected and can be positive, negative, or
    both.}
  \note{ISO 31000 counts positive deviations (opportunities) as risk. The EU
    AI Act and NIST AI RMF concentrate on harm. The EU AI Act classifies
    systems into risk tiers (prohibited, high-risk, limited, and minimal),
    and each tier carries different obligations.}
  \related{ai-governance,safety,red-teaming}
\end{entry}

\begin{entry}{alignment}{Alignment}[AI alignment]
  \defn{ji-2023-alignment}{The goal of making AI systems behave in line with
    human intentions and values. It is characterised by the RICE objectives:
    robustness, interpretability, controllability, and ethicality.}
  \defn{gabriel-2020}{The normative and technical problem of making AI act
    in accordance with appropriate values. The problem includes choosing
    whether a system should follow instructions, intentions, revealed
    preferences, ideal preferences, interests, or values.}
  \defn{russell-2019}{Ensuring that machines pursue objectives that are
    beneficial to humans. The proposed route is for machines to be uncertain
    about human preferences and to learn them from behaviour, rather than to
    optimise a fixed, possibly misspecified objective.}
  \related{rlhf,safety,human-oversight}
\end{entry}

\begin{entry}{auditability}{Auditability}[algorithmic audit]
  \defn{hleg-2019}{The ability to assess an AI system's algorithms, data, and
    design processes, by internal or external auditors, with the resulting
    reports available. It is especially important for systems affecting
    fundamental rights.}
  \defn{raji-2020}{Supported by an internal audit process, SMACTR (scoping,
    mapping, artifact collection, testing, and reflection), that produces
    documentation throughout development. The documentation lets an
    organisation check a system against its declared ethical expectations
    before deployment.}
  \defn{iso-42001}{Achieved through internal audits of the AI management
    system at planned intervals, which check conformity with the
    organisation's requirements and the standard.}
  \related{accountability,transparency,model-card}
\end{entry}

\begin{entry}{bias}{Bias}[AI bias]
  \defn{nist-sp-1270}{Bias in AI falls into three categories.
    \emph{Systemic} biases come from institutional practices.
    \emph{Statistical and computational} biases come from unrepresentative
    data or modelling choices. \emph{Human-cognitive} biases shape how
    people design, interpret, and use AI systems.}
  \defn{friedman-nissenbaum-1996}{Computer systems that \emph{systematically
    and unfairly discriminate} against certain individuals or groups in
    favour of others. Bias may be pre-existing, technical, or emergent.}
  \defn{mehrabi-2021}{Systematic errors or skews at any stage of the ML
    pipeline, such as data collection, sampling, measurement, algorithm
    design, or user interaction, that can lead to unfair outcomes.}
  \defn{hastie-2009}{In statistics, the difference between an estimator's
    expected value and the true value. This is a technical property with no
    moral meaning, and is often traded off against variance.}
  \note{The statistical meaning is neutral and often intentional, as in
    regularisation. The socio-technical meanings are about unjust outcomes.
    Say which meaning is intended.}
  \related{fairness,datasets}
\end{entry}

\begin{entry}{data-poisoning}{Data Poisoning}
  \defn{nist-ai-100-2}{An attack at training time in which an adversary
    controls part of the training data, or the labels, to degrade model
    performance in general (availability poisoning) or to cause targeted or
    backdoor behaviour (integrity poisoning).}
  \defn{biggio-2012}{Injecting specially crafted training points to raise a
    learning algorithm's test error. The paper demonstrates this against
    support vector machines by computing the points with gradient ascent.}
  \defn{owasp-llm-2025}{For large language models: manipulation of
    pre-training, fine-tuning, or embedding data to introduce
    vulnerabilities, backdoors, or biases that compromise security,
    performance, or ethical behaviour.}
  \related{adversarial-example,security-resilience,datasets}
\end{entry}

\begin{entry}{differential-privacy}{Differential Privacy}[DP]
  \defn{dwork-roth-2014}{A mathematical guarantee that the output of an
    analysis is almost equally likely whether or not any single individual's
    data are included. A mechanism $M$ is $(\varepsilon,\delta)$-differentially
    private if, for all neighbouring datasets $D, D'$ and all sets of outputs
    $S$, $\Pr[M(D)\in S] \le e^{\varepsilon}\Pr[M(D')\in S] + \delta$.}
  \defn[§3.6]{nist-ai-rmf}{One of several privacy-enhancing technologies
    that can support privacy-enhanced AI. It may trade off against accuracy,
    and therefore against fairness for small groups.}
  \related{privacy}
\end{entry}

\begin{entry}{explainability}{Explainability}[XAI]
  \defn[§3.5]{nist-ai-rmf}{A representation of the mechanisms underlying an
    AI system's operation. It answers the question of \emph{how} a decision
    was made.}
  \defn{nist-ir-8312}{Characterised by four principles. A system gives
    evidence or reasons for its outputs (\emph{explanation}). The
    explanations are understandable to their audience (\emph{meaningful}).
    They correctly reflect the process that produced the output
    (\emph{explanation accuracy}). The system operates only within its
    designed conditions and confidence (\emph{knowledge limits}).}
  \defn{barredo-2020}{Given an audience, an explainable AI is one that
    produces details or reasons to make its functioning clear or easy to
    understand.}
  \defn{hleg-2019}{Part of transparency: the ability to explain both the
    technical processes of an AI system and the related human decisions,
    adapted to the expertise of the stakeholder concerned.}
  \note{Every definition makes explainability depend on its \emph{audience}.
    NIST separates it from interpretability, but many papers use the two
    words interchangeably.}
  \related{interpretability,transparency}
\end{entry}

\begin{entry}{fairness}{Fairness}[algorithmic fairness]
  \defn[§3.7]{nist-ai-rmf}{The characteristic \enquote{fair --- with harmful
    bias managed}: concerns for equality and equity, addressed by managing
    harmful bias and discrimination. Standards of fairness are complex and
    differ across cultures and applications.}
  \defn{dwork-2012}{\emph{Individual fairness}: similar individuals should be
    treated similarly, formalised as a Lipschitz condition on the classifier
    under a task-specific similarity metric.}
  \defn{hardt-2016}{\emph{Equalised odds}: a predictor is fair if its
    predictions are independent of a protected attribute conditional on the
    true outcome. That is, true- and false-positive rates are equal across
    groups. \emph{Equal opportunity} relaxes this to equal true-positive
    rates.}
  \defn{hleg-2019}{A requirement named \emph{diversity, non-discrimination
    and fairness}. It covers avoiding unfair bias, accessibility and
    universal design, and involving stakeholders throughout the lifecycle.}
  \note{Group criteria, such as demographic parity and equalised odds, and
    individual criteria cannot in general all be satisfied at once. Choosing
    among them is a normative decision that should be documented.}
  \related{bias,accountability}
\end{entry}

\begin{entry}{human-oversight}{Human Oversight}
  \defn[Art.~14]{eu-ai-act}{High-risk AI systems must be designed and
    developed so that natural persons can effectively oversee them while
    they are in use. The aim is to prevent or minimise risks to health,
    safety, or fundamental rights, including by being able to understand,
    monitor, override, or stop the system.}
  \defn{hleg-2019}{Part of the requirement \emph{human agency and
    oversight}. It may be achieved through human-in-the-loop,
    human-on-the-loop, or human-in-command approaches, which differ in how
    closely a human is involved.}
  \defn{nist-ai-rmf}{Human roles and responsibilities in decision-making and
    in overseeing AI systems should be clearly defined and differentiated,
    taking account of human cognitive biases such as automation bias.}
  \note{Oversight is only effective if the overseer has the competence,
    authority, time, and information to intervene. Without these, a human
    reviewer may end up bearing responsibility for outcomes they could not
    actually control.}
  \related{human-in-the-loop,autonomy,accountability}
\end{entry}

\begin{entry}{human-in-the-loop}{Human-in-the-Loop}[HITL]
  \defn{hleg-2019}{The capability for human intervention in every decision
    cycle of the system. It contrasts with \emph{human-on-the-loop}, where
    humans intervene during design and monitor operation, and
    \emph{human-in-command}, where humans oversee overall activity and
    decide when and how to use the system.}
  \defn{mosqueira-rey-2023}{In machine learning, approaches in which humans
    take part in the learning process, for example through active learning,
    interactive ML, or machine teaching, so as to improve model accuracy or
    efficiency.}
  \note{The governance meaning concerns human control over decisions. The
    ML meaning concerns humans contributing to training.}
  \related{human-oversight,rlhf}
\end{entry}

\begin{entry}{interpretability}{Interpretability}
  \defn*{doshi-velez-2017}{The ability to explain or to present in
    understandable terms to a human.}
  \defn{miller-2019}{The degree to which a human can understand the cause of
    a decision. Miller uses it interchangeably with \emph{explainability}.}
  \defn[§3.5]{nist-ai-rmf}{The meaning of an AI system's output in the
    context of its designed functional purpose. It answers the question of
    \emph{why} a decision was made and what it means to the user.}
  \related{explainability,transparency}
\end{entry}

\begin{entry}{model-card}{Model Card}[datasheet]
  \defn{mitchell-2019}{A short document accompanying a trained model that
    gives benchmarked evaluation across relevant conditions, such as
    demographic groups, as well as intended use, out-of-scope uses,
    evaluation procedures, and ethical considerations.}
  \defn{gebru-2021}{The companion practice for data, called a
    \emph{datasheet for datasets}. It documents a dataset's motivation,
    composition, collection process, preprocessing, uses, distribution, and
    maintenance.}
  \defn[Art.~11]{eu-ai-act}{Related regulatory concept: the \emph{technical
    documentation} that providers of high-risk AI systems must draw up
    before placing them on the market and keep up to date.}
  \related{transparency,auditability,datasets}
\end{entry}

\begin{entry}{privacy}{Privacy}[privacy-enhanced AI]
  \defn[§3.6]{nist-ai-rmf}{The norms and practices that help to safeguard
    human autonomy, identity, and dignity. These include freedom from
    intrusion, limits on observation, and individuals' agency to consent to
    disclosure or control of facets of their identities.}
  \defn{hleg-2019}{The requirement \emph{privacy and data governance}: respect
    for privacy, the quality and integrity of data, and legitimate access to
    data, throughout the system's lifecycle.}
  \defn{nist-ai-100-2}{The target of \emph{privacy attacks}, such as
    membership inference, data reconstruction, and model extraction, which
    aim to recover information about training data or the model.}
  \related{differential-privacy,security-resilience}
\end{entry}

\begin{entry}{prompt-injection}{Prompt Injection}
  \defn{nist-ai-100-2}{An attack in which an adversary inserts instructions
    into a generative AI system's input. In \emph{direct} prompt injection
    the attacker is the user. In \emph{indirect} prompt injection the
    instructions arrive through external content, such as a web page or
    document, that the system consumes.}
  \defn{greshake-2023}{Indirect prompt injection: planting prompts in data
    likely to be retrieved by an LLM-integrated application, which blurs the
    line between data and instructions and allows the attacker to control
    the application remotely.}
  \defn{owasp-llm-2025}{A vulnerability, ranked first among risks to LLM
    applications, in which user prompts or external content alter the
    model's behaviour or output in unintended ways. Outcomes include data
    disclosure and unauthorised actions.}
  \note{Prompt injection is distinct from \emph{jailbreaking}, which targets
    the model's safety training. Prompt injection targets the application
    built around the model.}
  \related{prompt,agentic-ai,retrieval-augmented-generation,security-resilience}
\end{entry}

\begin{entry}{red-teaming}{Red Teaming}
  \defn{ganguli-2022}{Adversarially probing a language model, by human or
    automated means, to discover and measure harmful outputs, so that the
    harms can be analysed and reduced.}
  \defn{nist-ai-600-1}{A structured testing effort, often by teams outside
    the development group, to find flaws and vulnerabilities in a generative
    AI system, such as harmful or discriminatory outputs, unforeseen
    behaviour, or potential misuse.}
  \defn[Art.~55]{eu-ai-act}{Adversarial testing that providers of
    general-purpose AI models with systemic risk must carry out and
    document, in order to identify and mitigate systemic risks.}
  \related{ai-risk,safety,security-resilience}
\end{entry}

\begin{entry}{robustness}{Robustness}
  \defn{iso-22989}{The ability of a system to maintain its level of
    performance under any circumstances, including external interference
    and harsh environmental conditions.}
  \defn{hleg-2019}{Named \emph{technical robustness and safety}: AI systems
    are resilient to attack, have fall-back plans, and are accurate,
    reliable, and reproducible, so that unintentional harm is minimised.}
  \defn[Art.~15]{eu-ai-act}{High-risk AI systems shall be as resilient as
    possible to errors, faults, or inconsistencies within the system or its
    environment, for example through technical redundancy and fail-safe
    plans.}
  \defn{madry-2018}{\emph{Adversarial robustness}: low worst-case loss under
    any perturbation within an allowed set, formalised as a min–max
    (saddle-point) optimisation problem.}
  \related{adversarial-example,validity-reliability,distribution-shift}
\end{entry}

\begin{entry}{safety}{Safety}[AI safety]
  \defn[§3.2]{nist-ai-rmf}{AI systems should not, under defined conditions,
    lead to a state in which human life, health, property, or the
    environment is endangered.}
  \defn{amodei-2016}{The problem of \emph{accidents} in ML systems:
    unintended and harmful behaviour arising from poor design. Examples are
    negative side effects, reward hacking, unsafe exploration, and
    brittleness under distributional shift.}
  \defn{hleg-2019}{Part of technical robustness: AI systems are developed
    with a preventive approach to risk, behave reliably as intended, and
    minimise unintentional and unexpected harm.}
  \note{\enquote{AI safety} also names a research field concerned with
    risks from advanced AI, including misalignment. In product and
    engineering contexts it usually means freedom from unacceptable risk of
    physical or psychological harm.}
  \related{alignment,ai-risk,robustness}
\end{entry}

\begin{entry}{security-resilience}{Security and Resilience}
  \defn[§3.3]{nist-ai-rmf}{An AI system is \emph{resilient} if it can
    withstand unexpected adverse events or changes in its environment or use,
    or degrade safely and gracefully. It is \emph{secure} if it maintains
    confidentiality, integrity, and availability through mechanisms that
    prevent unauthorised access and use.}
  \defn[Art.~15]{eu-ai-act}{Cybersecurity: high-risk AI systems shall be
    resilient against attempts by unauthorised third parties to alter their
    use, outputs, or performance by exploiting vulnerabilities. Examples
    include data poisoning, model poisoning, adversarial examples, and
    confidentiality attacks.}
  \defn{nist-ai-100-2}{Adversarial machine learning classifies attacks by
    the attacker's goals (availability, integrity, privacy, and misuse),
    capabilities, knowledge, and the stage of the ML lifecycle they target.}
  \related{adversarial-example,data-poisoning,prompt-injection,robustness}
\end{entry}

\begin{entry}{transparency}{Transparency}
  \defn[§3.4]{nist-ai-rmf}{The extent to which information about an AI
    system and its outputs is available to the individuals who interact with
    it, whether or not they know they are doing so.}
  \defn{hleg-2019}{A requirement covering \emph{traceability} of data and
    processes, \emph{explainability}, and \emph{communication}. People are
    told when they are interacting with an AI system, and they are told its
    capabilities and limitations.}
  \defn[Art.~13]{eu-ai-act}{High-risk AI systems shall be designed so that
    their operation is sufficiently transparent to enable deployers to
    interpret the system's output and use it appropriately. The systems are
    accompanied by instructions for use.}
  \defn[Art.~50]{eu-ai-act}{Transparency obligations toward the people
    affected: for example, disclosing that they are interacting with an AI
    system, and marking synthetic or deep-fake content.}
  \related{explainability,accountability,model-card}
\end{entry}

\begin{entry}{trustworthy-ai}{Trustworthy AI}
  \defn{hleg-2019}{AI that is \emph{lawful}, \emph{ethical}, and
    \emph{robust}. It is realised through seven requirements: human agency
    and oversight; technical robustness and safety; privacy and data
    governance; transparency; diversity, non-discrimination and fairness;
    societal and environmental well-being; and accountability.}
  \defn[§3]{nist-ai-rmf}{AI that has these characteristics: valid and
    reliable; safe; secure and resilient; accountable and transparent;
    explainable and interpretable; privacy-enhanced; and fair with harmful
    bias managed. Validity and reliability are the base for the others, and
    the characteristics must be balanced against each other in context.}
  \defn{iso-24028}{Rests on \emph{trustworthiness}: the ability to meet
    stakeholders' expectations in a verifiable way. For AI this covers
    characteristics such as reliability, availability, resilience, security,
    privacy, safety, accountability, transparency, integrity, authenticity,
    quality, and usability.}
  \defn{oecd-2024}{AI that respects the OECD values-based principles:
    inclusive growth, sustainable development and well-being; human rights
    and democratic values, including fairness and privacy; transparency and
    explainability; robustness, security and safety; and accountability.}
  \note{The frameworks overlap heavily but are organised differently. The
    EU HLEG starts from ethics and fundamental rights, NIST from risk
    management, and ISO/IEC from verifiable stakeholder expectations.}
  \related{accountability,explainability,fairness,privacy,robustness,safety,transparency,validity-reliability}
\end{entry}

\begin{entry}{validity-reliability}{Validity and Reliability}
  \defn[§3.1]{nist-ai-rmf}{\emph{Validation} is confirmation, through
    objective evidence, that the requirements for a specific intended use
    have been fulfilled. \emph{Reliability} is the ability of an item to
    perform as required, without failure, for a given time interval under
    given conditions. Accuracy and robustness contribute to both.}
  \defn{iso-22989}{Reliability is the property of consistent intended
    behaviour and results.}
  \defn{hleg-2019}{Part of technical robustness: a reliable system works
    properly across a range of inputs and situations, and reproducibility
    means it behaves the same way when an experiment is repeated under the
    same conditions.}
  \note{NIST treats validity and reliability as the foundation of
    trustworthiness. A system that is not valid for its context cannot be
    made trustworthy by adding the other characteristics.}
  \related{robustness,trustworthy-ai,distribution-shift}
\end{entry}
